\subsection{Dimension}
\paragraph*{Defintion}
If $U$ is a subspace of $\mathbb{R}^n$, then a set $B=\{\vec{x}_1, \cdots, \vec{x}_k\} \subseteq \mathbb{R}^n$ is called a \textbf{basis} for $U$ if 
\begin{enumerate}
    \item $B$ is linearly independent, and
    \item $B$ spans $U$.
\end{enumerate}

\paragraph*{Theorem}
If $\{\vec{x}_1, \cdots, \vec{x}_m\}$ and $\{\vec{y}_1, \cdots, \vec{y}_p\}$ are both bases for a subspace $U$ of $\mathbb{R}^n$, then $m=p$.

\paragraph*{Definition}
If $\{\vec{x}_1, \cdots, \vec{x}_k\}$ is a basis for a subspace $U$ of $\mathbb{R}^n$, then the \textbf{dimension} of $U$, denoted $\dim(U)$, is defined to be $k$, the number of vectors in the basis.

\paragraph*{Example}
$\dim{\mathbb{R}^n}=n \implies$ a basis of $\mathbb{R}^n$ has $n$ vectors, for example $\{\vec{e}_1, \cdots, \vec{e}_n\}$, the standard unit vectors in $\mathbb{R}^n$.

\paragraph*{Properties}
\begin{itemize}
    \item You need \emph{at least} $n$ vectors to span $\mathbb{R}^n$.
    \item You can have \emph{at most} $n$ linearly independent vectors in $\mathbb{R}^n$.
    \item Any set of $n$ linearly independent vectors in $\mathbb{R}^n$ forms a basis for $\mathbb{R}^n$.
    \item Any set of $n$ vectors that spans $\mathbb{R}^n$ forms a basis for $\mathbb{R}^n$.
\end{itemize}

\paragraph*{Example}
Let $U = \{\begin{bmatrix}
    r \\
    s \\
    r
\end{bmatrix} \mid r, s \in \mathbb{R}\}$. Show that $U$ is a subspace of $\mathbb{R}^3$ and find a basis and the dimension of $U$.
\begin{align*}
    \begin{bmatrix}
        r \\
        s \\
        r
    \end{bmatrix} &= r \begin{bmatrix}
        1 \\
        0 \\
        1
    \end{bmatrix} + s \begin{bmatrix}
        0 \\
        1 \\
        0
    \end{bmatrix} \implies U = \text{span} \left\{ \begin{bmatrix}
        1 \\
        0 \\
        1
    \end{bmatrix}, \begin{bmatrix}
        0 \\
        1 \\
        0
    \end{bmatrix} \right\} \\
    \implies &U \text{ has a basis } \left\{ \begin{bmatrix}
        1 \\
        0 \\
        1
    \end{bmatrix}, \begin{bmatrix}
        0 \\
        1 \\
        0
    \end{bmatrix} \right\} \\
    &\implies \dim(U) = 2  
\end{align*}

\paragraph*{Example}
Let $V = \left\{ \begin{bmatrix}
    a + b \\
    2a - c \\
    c + 3a \\
    a - b
\end{bmatrix} \mid a, b, c \in \mathbb{R} \right\}$. Is $V$ a subspace of $\mathbb{R}^4$? If so, find a basis and the dimension of $V$.
\begin{align*}
    \implies a \begin{bmatrix}
        1 \\
        2 \\
        3 \\
        1
    \end{bmatrix} + b \begin{bmatrix}
        1 \\
        0 \\
        0 \\
        -1
    \end{bmatrix} + c \begin{bmatrix}
        0 \\
        -1 \\
        1 \\
        0
    \end{bmatrix} \implies B = \left\{ \begin{bmatrix}
        1 \\
        2 \\
        3 \\
        1
    \end{bmatrix}, \begin{bmatrix}
        1 \\
        0 \\
        0 \\
        -1
    \end{bmatrix}, \begin{bmatrix}
        0 \\
        -1 \\
        1 \\
        0
    \end{bmatrix} \right\} \text{ spans } V \text{ and all elements of } B \text{ are linearly independent, so } B \text{ is a basis for } V \text{ and } \dim(V) = 3.
\end{align*}

\paragraph*{Theorem}
Any subspace $U \neq \{\vec{0}\}$ has a basis and the dimension of $U$ is constant.

\paragraph*{Theorem}
Let $B = \{\vec{x}_1, \cdots, \vec{x}_k\} \subseteq \mathbb{R}^n$ be a basis for a subspace $U$.
If $A_{n \times n}$ is invertible, then $B_A = \{A\vec{x}_1, \cdots, A\vec{x}_k\}$ is also a basis for $U_A = \{A\vec{u} \mid \vec{u} \in U\}$.

\paragraph*{Theorems/Facts}
If $U \neq \{\vec{0}\}$ is a subspace of $\mathbb{R}^n$, then:
\begin{enumerate}
    \item $U$ has a basis and $\dim(U) \leq n$.
    \item Any independent set $B$ in $U$ can be extended to a basis for $U$ by adding vectors that are independent of elements of $B$. We can use elements from any fixed basis for $U$ to extend $B$ to a basis.
    \item Any spanning set $B$ for $U$ can be cut down to a basis of $U$ by removing vectors that are linear combinations of the others in $B$ until the remaining set is linearly independent and still spans $U$.
    \item Check for linear independence by finding determinant and checking that it is nonzero (for a square matrix formed from the vectors as columns).
\end{enumerate}

\paragraph*{Example}
Find a basis for $\mathbb{R}^4$ containing $B = \{\vec{u}, \vec{v}\}$ where $\vec{u} = \begin{bmatrix}
    0 \\
    1 \\
    2 \\
    3
\end{bmatrix}, \quad \vec{v} = \begin{bmatrix}
    2 \\
    -1 \\
    0 \\
    1
\end{bmatrix}$
We can extend $B$ by adding $\vec{e}_1$ and $\vec{e}_2$ to get a basis for $\mathbb{R}^4$:
\[
B' = \left\{ \vec{u}, \vec{v}, \vec{e}_1, \vec{e}_2 \right\} = \left\{ \begin{bmatrix}
    0 \\
    1 \\
    2 \\
    3
\end{bmatrix}, \begin{bmatrix}
    2 \\
    -1 \\
    0 \\
    1
\end{bmatrix}, \begin{bmatrix}
    1 \\
    0 \\
    0 \\
    0
\end{bmatrix}, \begin{bmatrix}
    0 \\
    1 \\
    0 \\
    0
\end{bmatrix} \right\}.
\]

\paragraph*{Example}
Let $A = \begin{bmatrix}
      1 & 3 & -1 & 7 \\
      2 & 6 & -2 & 14
\end{bmatrix}$.
\begin{align*}
    \text{IM}(A) = \text{span} \left\{ \begin{bmatrix}
        1 \\
        2
    \end{bmatrix}, \begin{bmatrix}
        3 \\
        6
    \end{bmatrix}, \begin{bmatrix}
        -1 \\
        -2
    \end{bmatrix}, \begin{bmatrix}
        7 \\
        14
    \end{bmatrix} \right\} \\
    \text{RREF}(A) = \begin{bmatrix}
        1 & 3 & -1 & 7 \\
        0 & 0 & 0 & 0
    \end{bmatrix} \implies \text{IM}(A) = \text{span} \left\{ \begin{bmatrix}
        1 \\
        2
    \end{bmatrix} \right\} \text{ and } \dim(\text{IM}(A)) = 1
\end{align*}

\paragraph*{Note}
When finding the basis for the image of a vector, reduce it to RREF and take the pivot columns of the original matrix as the vectors that form a basis for the image.